% SGA 3, Expose VI_B, Loop 1 English reconstruction.
% Scope: Proposition 1.4 through Corollary 1.4.2 in full;
% Counterexample 1.4.3 is excluded.
% Authority: Polo--Gille Exp6B-13oct24.pdf, component-local pages 5--7
% (combined-reader pages 346--348; original-margin printed pages 325--327).
% OCR and the Reinhold English Markdown are comparison witnesses only.

\medskip
\phantomsection\label{prop:sga3-VIB-1.4}%
\noindent\textbf{命题 1.4.---}
\emph{设 \(A\) 为 Artin 局部环, \(u:G\to H\)
为局部有限型的 \(A\)-群之间的拟紧态射. 以下断言等价:
\begin{enumerate}
\item[\textup{\SGAInlineRefTarget{prop:sga3-VIB-1.4:item:i}(i)}] \(u\) 固有;
\item[\textup{\SGAInlineRefTarget{prop:sga3-VIB-1.4:item:ii}(ii)}]
  存在 \(h\in H\), 使纤维 \(u^{-1}(h)\) 非空且在 \(\kappa(h)\) 上固有;
\item[\textup{\SGAInlineRefTarget{prop:sga3-VIB-1.4:item:iii}(iii)}]
  \(\ker u\) 在 \(A\) 上固有.
\end{enumerate}}

\emph{证明.}
显然, \textup{\SGARefLink{prop:sga3-VIB-1.4:item:i}{(i)}} 蕴含
\textup{\SGARefLink{prop:sga3-VIB-1.4:item:iii}{(iii)}}, 而
\textup{\SGARefLink{prop:sga3-VIB-1.4:item:iii}{(iii)}} 蕴含
\textup{\SGARefLink{prop:sga3-VIB-1.4:item:ii}{(ii)}}.
假设蕴含 \(u\) 为有限型; 又由于 \(G\) 分离 (第~VI\(_A\)~篇, 0.3),
\(u\) 也分离 (EGA~I, 5.5.1).
只须证明 \textup{(ii)} 蕴含 \(u\) 普遍闭.
因而可假设 \(A\) 等于其剩余域 \(k\).%
\footnote{\label{note:sga3-VIB-12}编者注: 添加了下面一句.}
令 \(k'\) 为 \(\kappa(h)\) 的一个代数闭包, 设
\[
  u':G'\longrightarrow H'
\]
由 \(u\) 经基变换得到, 并令 \(h'\in H'\) 为 \(h\) 上方的点. 则
\[
  u'^{-1}(h')=u^{-1}(h)\times_{\kappa(h)}k'
\]
非空且固有, 只须证明 \(u'\) 固有即可 (EGA~IV\(_2\), 2.6.4).
因此可假设 \(k\) 代数闭且 \(h\in H(k)\).

由 \hyperref[prop:sga3-VIB-1.2]{1.2}, 此时 \(u(G)\) 是
\(H\) 的一个约化闭群子概形的底集.
由于每个闭浸入都固有 (EGA~II, 5.4.2), 可假设 \(u\) 满射且 \(H\) 约化.
此时群 \(G(k)\) 在 \(H\) 的闭点集上传递地作用.
因此, 对每个闭点 \(y\in H\), 纤维 \(u^{-1}(y)\)
在 \(\kappa(y)\) 上固有. 由 EGA~IV\(_3\), 9.6.1,
使 \(u^{-1}(y)\) 在 \(\kappa(y)\) 上不固有的 \(y\in H\)
组成的集合局部可构造. 它不包含闭点, 因而为空
(参见 EGA~IV\(_3\), 10.3.1 和 10.4.7).

\footnote{\label{note:sga3-VIB-13}编者注: 以下展开了原文.}
现考虑 \(H^0\) 的泛点 \(\eta\). 由上面的论证, 纤维
\[
  u^{-1}(\eta)=G\times_H\operatorname{Spec}\kappa(\eta)
\]
在 \(\kappa(\eta)\) 上固有.
由于 \(H\) 约化, \(\kappa(\eta)=\mathcal O_{H,\eta}\).
后一个环是环 \(\mathcal O_H(V)\) 的正向极限,
其中 \(V\) 遍历 \(H^0\) 的非空开子集.
由 EGA~IV\(_3\), 8.1.2(a) 和 8.10.5(xii),
存在非空开子集 \(V\subset H^0\), 使 \(u\) 在 \(V\) 上的限制固有.
当 \(g\in G(k)\) 时, 各平移 \(g\cdot V\) 构成 \(H\) 的开覆盖;
且对每个 \(g\in G(k)\), \(u\) 在 \(g\cdot V\) 上的限制固有.
故 \(u\) 固有 (参见 EGA~II, 5.4.1).

\medskip
\phantomsection\label{cor:sga3-VIB-1.4.1}%
\noindent\textbf{推论 1.4.1.---}
\emph{设 \(A\) 为 Artin 局部环, \(u:G\to H\)
为局部有限型的 \(A\)-群之间的态射. 以下断言等价:
\begin{enumerate}
\item[\textup{(i)}] \(u\) 局部拟有限;
\item[\textup{(ii)}] \(u\) 在某一点拟有限;
\item[\textup{(iii)}] \(\ker u\) 离散;
\item[\textup{\SGAInlineRefTarget{cor:sga3-VIB-1.4.1:item:iv}(iv)}]
  \(u\) 在 \(G\) 的每个连通分支上的限制有限.
\end{enumerate}
若 \(u\) 拟紧, 则这些断言还等价于
\begin{enumerate}
\item[\textup{(v)}] \(u\) 有限.
\end{enumerate}}

显然, \textup{(iv)} 蕴含 \textup{(iii)},
\textup{(iii)} 蕴含 \textup{(ii)} (EGA~I, 6.4.4),
且当 \(u\) 拟紧时, \textup{(iv)} 与 \textup{(v)} 等价.
在 \hyperref[prop:sga3-VIB-1.3]{1.3} 中已经看到,
\textup{\SGARefLink{concl:sga3-VIB-1.3-i}{(i)}} 与
\textup{\SGARefLink{prop:sga3-VIB-1.3:item:ii}{(ii)}} 等价.

只须证明 \textup{\SGARefLink{concl:sga3-VIB-1.3-i}{(i)}} 蕴含
\textup{\SGARefLink{cor:sga3-VIB-1.4.1:item:iv}{(iv)}}.
设 \(C\) 为 \(G\) 的一个连通分支.
由于 \(C\) 在 \(A\) 上为有限型
(第~VI\(_A\)~篇, \SGARefLink{cor:sga3-VIA-2.4.1}{2.4.1}),
且 \(G,H\) 分离 (第~VI\(_A\)~篇, \SGARefLink{sga3:VIA:0.3}{0.3}),%
\begingroup
\renewcommand{\thefootnote}{\fnsymbol{footnote}}%
\footnote[1]{\label{note:sga3-VIB-1.4.1-separated-xref}译者注:
法文底本在此引用第~VI\(_A\)~篇, 0.2;
分离性陈述是 \SGARefLink{sga3:VIA:0.3}{0.3},
上面的编者\hyperref[note:sga3-VIB-2]{注 2} 也回顾了这一点.}%
\endgroup{}
由 EGA~I, 5.5.1 和 6.3.4 可知, \(u\) 在 \(C\)
上的限制 \(u'\) 分离且为有限型.%
\footnote{\label{note:sga3-VIB-14}编者注: 以下展开了原文.}
它的纤维离散, 故 \(u'\) 拟有限 (参见 EGA~II, 6.2.2).
由于每个拟有限的固有态射都有限 (参见 EGA~III\(_1\), 4.4.2),
只须证明 \(u'\) 普遍闭.

为此, 可假设 \(A\) 等于其剩余域 \(k\).
由 \(\mathrm{fpqc}\) 下降 (参见 EGA~IV\(_2\), 2.6.4),
只须证明
\[
  u'\otimes_k\overline k
\]
普遍闭, 其中 \(\overline k\) 是 \(k\) 的一个代数闭包.
由于 \(C\) 在 \(k\) 上为有限型,
概形 \(C\otimes_k\overline k\) 是 \(G\otimes_k\overline k\)
的有限个连通分支 \(C'_1,\ldots,C'_n\) 的不交并,
只须对每个 \(C'_i\) 证明该断言.
因而化归到 \(k\) 代数闭的情形.

设 \(g\) 为 \(C\) 的一个闭点, 令
\[
  u^0:G^0\longrightarrow H
\]
为 \(u\) 在 \(G^0\) 上的限制. 则
\[
  u'=r_{u(g)}\circ u^0\circ r_{g^{-1}},
\]
其中 \(r_g\) 表示以 \(g\) 作右平移.
因此, 为证明 \(u'\) 固有, 只须证明 \(u^0\) 固有.
由假设 \(u\) 局部拟有限, 故非空纤维 \(\ker u\) 离散.
我们已看到 \(u^0\) 为有限型, 因而其纤维 \(\ker u^0\) 有限
(参见 EGA~II, 6.2.2), 故非空且固有.
现在由 \hyperref[prop:sga3-VIB-1.4]{命题~1.4}, \(u^0\) 固有.

\medskip
\phantomsection\label{cor:sga3-VIB-1.4.2}%
\noindent\textbf{推论 1.4.2.---}
\emph{设 \(A\) 为 Artin 局部环, \(u:G\to H\)
为局部有限型的 \(A\)-群之间的拟紧态射. 以下断言等价:%
\footnote{\label{note:sga3-VIB-15}编者注: 可以去掉 \(G\) 和 \(H\)
局部有限型的假设, 因为 \textup{[Per76]}, 4.2.4 陈述:
域 \(k\) 上的群概形之间的每个拟紧单态射 \(u:G\to H\) 都是闭浸入.
当 \(G\) 和 \(H\) 仿射时, 另见 \textup{[DG70]}, III.3.7.2(b);
此时每个态射 \(G\to H\) 都仿射 (EGA~II, 1.6.2(v)), 因而拟紧.}
\begin{enumerate}
\item[\textup{\SGAInlineRefTarget{cor:sga3-VIB-1.4.2:item:i}(i)}] \(u\) 是闭浸入;
\item[\textup{\SGAInlineRefTarget{cor:sga3-VIB-1.4.2:item:ii}(ii)}] \(u\) 是单态射;
\item[\textup{\SGAInlineRefTarget{cor:sga3-VIB-1.4.2:item:iii}(iii)}]
  \(\ker u\) 平凡, 即同构于单位 \(A\)-群.
\end{enumerate}
特别地, \(H\) 的每个群子概形%
\footnote{\label{note:sga3-VIB-16}编者注: 在这一特殊情形中,
另见第~VI\(_A\)~篇, \SGARefLink{sga3:VIA:0.5.2}{0.5.2},
该结果无须任何有限性假设.}
都是闭的.}

\begingroup
\renewcommand{\thefootnote}{\fnsymbol{footnote}}%
\footnote[1]{\label{note:sga3-VIB-1.4.2-unit-base}译者注:
法文底本在陈述及证明中都将平凡核称为``单位 \(k\)-群''.
这里的基底是 Artin 局部环 \(A\), 而 \(\ker u\) 是 \(A\)-群;
因而单态射判别法中所需的单位对象应为 \(\operatorname{Spec}A\),
即单位 \(A\)-群.}%
\endgroup{}
显然, \textup{\SGARefLink{cor:sga3-VIB-1.4.2:item:i}{(i)}} 蕴含
\textup{\SGARefLink{cor:sga3-VIB-1.4.2:item:ii}{(ii)}}.
考虑 \(G\) 和 \(H\) 表示的函子, 立即可见
\textup{\SGARefLink{cor:sga3-VIB-1.4.2:item:ii}{(ii)}} 与
\textup{\SGARefLink{cor:sga3-VIB-1.4.2:item:iii}{(iii)}} 等价.
最后, 若 \(\ker u\) 为单位 \(A\)-群, 则它是非空的固有纤维.
因此 \hyperref[prop:sga3-VIB-1.4]{命题~1.4} 表明 \(u\) 为固有单态射;
由于 \(H\) 局部 Noether, 它为有限呈示 (EGA~IV\(_1\), 1.6.1),
因而为闭浸入 (EGA~IV\(_3\), 8.11.5).

最后一个断言来自 \(H\) 局部 Noether,
因而每个浸入 \(G\to H\) 都拟紧 (EGA~I, 6.6.4).
% Translation statistics: source_lines=194; translated_lines=180; inline_math=141; display_math=6; pending_TODO=0.
% Review: complete component translation; math fragments, reference keys, environments and note counts checked against source; formula prose translated; compilation pending parent integration.
